% TEMPLATE-MILC-v3.tex - plantilla maestra UNICA de las guias MILC v3.
%
% La usan las cuatro familias de guias, cada una con su driver:
%   · Grados 6-11          -> scripts/build-guias-g11.py
%   · Bebras               -> scripts/build-guias-bebras.py
%   · Semillero            -> scripts/build-guias-semillero.py
%   · Territorio interior  -> scripts/build-guias-territorio-interior.py
%
% Antes habia tres copias casi identicas (~400 lineas iguales, 6 distintas):
% cada mejora habia que aplicarla tres veces, y en la practica no se hacia
% (el motor de recursos vivia solo en dos de ellas). Lo que varia entre
% familias esta parametrizado en 6 placeholders que cada driver llena
% (se nombran aqui SIN su sintaxis real: el builder reemplaza por texto plano,
% tambien dentro de los comentarios, y un valor multilinea rompe el preambulo):
%
%   HEADER_IZQ        encabezado izquierdo de pagina
%   HEADER_DER        encabezado derecho de pagina
%   PORTADA_ETIQUETA  rotulo sobre el numero grande (GRADO/MOMENTO/MODULO)
%   PORTADA_SERIE     linea de serie de la portada
%   PORTADA_ID        identificador de la guia en la portada
%   PORTADA_CIRCULO   el circulito de grado (°), o vacio si no aplica
%   PDF_TITULO        titulo en texto plano (sin \textbf ni comillas TeX) para
%                     los metadatos del PDF (/Title); lo produce lib_xelatex.texto_pdf
%
% Composicion (auditoria 2026-09-03): las bandas de estacion piden espacio
% (\Needspace*) y prohiben el salto justo despues (\nopagebreak), asi no quedan
% huerfanas al pie; las cajas largas (softbox, drawbox, infoband, puentebox) son
% `breakable`, asi una caja de media pagina ya no arrastra todo a la siguiente
% dejando la anterior al 40 %. Todo texto sobre mostaza o verde va en milcNegro
% (blanco sobre #E5B400 da 1,93:1 y sobre #58A923 2,95:1; ninguno pasa WCAG AA).
% El resto de claves las documenta content/guias/_SCHEMA.md. El script valida
% que no queden placeholders sin reemplazar antes de escribir el archivo final.

% Etiquetado PDF (latex-lab): /Lang, estructura y texto alternativo de las
% figuras, para lectores de pantalla. Debe ir ANTES de \documentclass.
\DocumentMetadata{lang=es-CO,tagging=on}
\documentclass[11pt,letterpaper]{article}
% headheight=24pt: la cabecera derecha lleva dos lineas (identidad de la guia
% y estacion actual). headsep se acorta en la misma medida para que el bloque
% de texto quede exactamente donde estaba con headheight=12pt.
\usepackage[margin=2.0cm,headheight=24pt,headsep=13pt]{geometry}
\usepackage[table]{xcolor}
\usepackage{fontspec}
\usepackage[spanish]{babel}
\usepackage{tikz}
% Librerías TikZ para los diagramas de las guías (bloque `recursos:`):
%  · babel  → babel en español vuelve activos caracteres como «>», y TikZ los usa
%             en sus opciones (>=stealth, ->). Sin esta librería, cualquier
%             diagrama con flechas revienta con «\language@active@arg> has an
%             extra }». Debe ir ANTES de que se dibuje nada.
%  · positioning        → sintaxis «below left=2pt and 3pt».
%  · decorations.pathreplacing → llaves ({brace}) para señalar rangos.
\usetikzlibrary{babel,positioning,decorations.pathreplacing}
\usepackage{graphicx}
% Circuitos: el trabajo de electrónica se hace hoy con micro:bit y sus sensores
% integrados (MakeCode, sin cableado), así que no se necesitan esquemáticos.
% Si algún día entran montajes con protoboard, basta descomentar esta línea:
% los diagramas .tex/.tikz declarados en `recursos:` ya se incrustan solos.
% \usepackage{circuitikz}
\usepackage[most]{tcolorbox}
\usepackage{tabularx}
\usepackage{array}
\usepackage{fancyhdr}
\usepackage{titlesec}
\usepackage[document]{ragged2e}  % bandera a la derecha en todo el cuerpo
\usepackage{enumitem}
\usepackage{microtype}
\usepackage{etoolbox}
\usepackage{needspace}
% hyperref al final del preambulo: metadatos del PDF (titulo, autor, idioma,
% familia), marcadores por estacion (\pdfbookmark) y enlaces sin recuadro.
\usepackage[hidelinks,
  pdftitle={Variables y umbrales — calibrar antes de programar},
  pdfauthor={PhD. Álvaro Cárdenas Orozco},
  pdfsubject={Tecnología e Informática · Grado 8 · Guía 16 · MILC v3},
  pdflang={es-CO},
  bookmarksopen=true,bookmarksnumbered=false]{hyperref}

% Helvetica Neue no trae la flecha U+2192, asi que cada «→» escrito literalmente
% en el YAML se compilaba a NADA, en silencio (462 flechas en 61 guias: rutas del
% tipo «paleta → tipografia → lockup» salian rotas). Mapeamos el caracter a su
% version matematica, que si tiene glifo. Asi el autor puede seguir escribiendo
% «→» comodamente en el YAML.
\usepackage{newunicodechar}
\newunicodechar{→}{\ensuremath{\rightarrow}}
% Mismo problema con los simbolos de marca y vinetas: el «✓» de «una palomita ✓»
% salia como cuadrito vacio en el PDF de todas las guias que lo usaban.
\usepackage{pifont}
\newunicodechar{✓}{\ding{51}}
\newunicodechar{✗}{\ding{55}}
\newunicodechar{✔}{\ding{52}}
\newunicodechar{•}{\textbullet}
\newunicodechar{≥}{\ensuremath{\geq}}
\newunicodechar{≤}{\ensuremath{\leq}}

\setmainfont{Helvetica Neue}
\setsansfont{Helvetica Neue}
\setlength{\parindent}{0pt}
\setlength{\parskip}{6pt}
% Sin particion de palabras: en 6.º-7.º una palabra cortada por guion al final
% de linea («Comuni-cacion») frena la lectura mucho mas que una linea corta.
\hyphenpenalty=10000
\exhyphenpenalty=10000
\pretolerance=2000
\tolerance=3000
\setlength{\emergencystretch}{3em}
\linespread{1.10}
\renewcommand{\arraystretch}{1.25}
\setlist{leftmargin=5mm,itemsep=1mm,topsep=1mm}
\newcolumntype{Y}{>{\RaggedRight\arraybackslash}X}

\definecolor{milcMagenta}{HTML}{E6007E}
\definecolor{milcVino}{HTML}{5A0038}
\definecolor{milcTurquesa}{HTML}{0093A5}
\definecolor{milcVerde}{HTML}{58A923}
\definecolor{milcMostaza}{HTML}{E5B400}
\definecolor{milcOcre}{HTML}{B86600}
\definecolor{milcNegro}{HTML}{241816}
\definecolor{milcGris}{HTML}{F7F7F4}
\definecolor{milcLinea}{HTML}{E8E2DD}
\definecolor{milcGradePrimary}{HTML}{FF6600}
\definecolor{milcGradeDark}{HTML}{9A3E00}
\definecolor{milcGradeSoft}{HTML}{FFE4D0}
\definecolor{milcGradeAccent}{HTML}{CFE7FF}
\definecolor{milcGradeText}{HTML}{FFFFFF}
\color{milcNegro}

\pagestyle{fancy}
\fancyhf{}
% Cabecera de dos lineas: identidad de la guia arriba y, a la derecha y en
% gris, la estacion en curso (\markright desde \milcestacion). Antes de la
% primera estacion la segunda linea va vacia; el \strut mantiene la altura
% para que la linea de arriba no baile entre paginas.
\lhead{\small Tecnología e Informática · Grado 8\\[-1pt]{\footnotesize\strut}}
\rhead{\small Guía 16 · MILC v3\\[-1pt]{\footnotesize\strut\color{milcNegro!65}\rightmark}}
\cfoot{\small\thepage}
\renewcommand{\headrulewidth}{0pt}

% Sobre mostaza (#E5B400) y verde (#58A923) el texto va en milcNegro; sobre el
% resto de la paleta, en blanco. Se usa dentro de fonttitle/fontupper, donde un
% \color posterior manda sobre coltitle.
\newcommand{\milctintasobre}[1]{%
  \ifboolexpr{test{\ifstrequal{#1}{milcMostaza}} or test{\ifstrequal{#1}{milcVerde}}}%
    {\color{milcNegro}}{\color{white}}}

\newtcolorbox{titlebox}[1]{
  enhanced,colback=#1,colframe=#1,arc=7mm,boxrule=0pt,
  left=7mm,right=7mm,top=3mm,bottom=3mm,
  before skip=8pt,after skip=8pt,
  fontupper=\sffamily\bfseries\Large\color{white}
}

\newtcolorbox{softbox}[3]{
  enhanced,breakable,pad at break=3mm,
  colback=#2,colframe=#1,borderline west={4pt}{0pt}{#1},
  arc=2mm,boxrule=.4pt,left=4mm,right=4mm,top=3mm,bottom=3mm,
  before skip=6pt,after skip=8pt,title={#3},coltitle=white,
  colbacktitle=#1,fonttitle=\sffamily\bfseries\milctintasobre{#1},fontupper=\RaggedRight,
  attach boxed title to top left={xshift=3mm,yshift=-2mm},
  boxed title style={arc=2mm, boxrule=0pt}
}

\newtcolorbox{drawbox}[2]{
  enhanced,breakable,pad at break=4mm,
  colback=white,colframe=#1,arc=4mm,boxrule=1pt,
  left=5mm,right=5mm,top=4mm,bottom=4mm,before skip=8pt,after skip=8pt,
  title={#2},coltitle=white,colbacktitle=#1,fonttitle=\sffamily\bfseries\milctintasobre{#1},
  fontupper=\RaggedRight,
  attach boxed title to top left={xshift=4mm,yshift=-2mm},
  boxed title style={arc=3mm, boxrule=0pt}
}

\newtcolorbox{infoband}[1]{
  enhanced,breakable,pad at break=2mm,colback=#1!8,colframe=#1!50,arc=2mm,boxrule=.5pt,
  left=4mm,right=4mm,top=2mm,bottom=2mm,before skip=4pt,after skip=4pt,
  fontupper=\small\RaggedRight
}

% Puente narrativo entre secciones (contrato editorial Conéctate).
% Da continuidad: "Acabas de X. Ahora vas a Y porque Z."
% Visualmente: banda izquierda en color del grado, texto en itálica pequeña.
\newtcolorbox{puentebox}{
  enhanced,breakable,pad at break=2mm,colback=milcGradeSoft,colframe=milcGradeDark,
  borderline west={3pt}{0pt}{milcGradeDark},
  arc=0mm,boxrule=0pt,left=5mm,right=4mm,top=2.5mm,bottom=2.5mm,
  before skip=4pt,after skip=8pt,
  fontupper=\itshape\small\color{milcNegro}\RaggedRight
}
% La flecha va en modo matematico a proposito: Helvetica Neue NO tiene el glifo
% U+2192, asi que \textrightarrow se compilaba a NADA («Missing character: There
% is no → in font Helvetica Neue») y el puente salia sin flecha. En modo
% matematico la flecha viene de la fuente matematica y si se dibuja.
\newcommand{\puente}[1]{%
  \begin{puentebox}%
    \textcolor{milcGradeDark}{$\rightarrow$}\hspace{2mm}#1%
  \end{puentebox}%
}

\newcommand{\linead}[1]{\noindent\color{milcLinea}\rule{#1}{0.5pt}\color{milcNegro}}
\newcommand{\respuesta}[1]{\par\vspace{2mm}\foreach \n in {1,...,#1}{\noindent\color{milcLinea}\rule{\linewidth}{0.5pt}\par\vspace{5mm}}\color{milcNegro}}
\newcommand{\checkbox}{\textcolor{milcGradeDark}{\fbox{\phantom{x}}}\hspace{5pt}}

% ─── Listas aireadas para las guías socioemocionales ────────────────────
% Componer las verificaciones como «\checkbox texto\\» deja el texto que salta
% de línea debajo del cuadrito y apelmaza el bloque. Estos entornos alinean la
% continuación y airean los ítems. Los usa Territorio Interior; cualquier
% familia puede hacerlo.
\newlist{checklista}{itemize}{1}
\setlist[checklista]{
  label=\textcolor{milcGradeDark}{\fbox{\phantom{x}}},
  leftmargin=9mm, labelsep=3.2mm, itemsep=2.4mm,
  topsep=2.5mm, parsep=0pt, partopsep=0pt, before=\RaggedRight
}
\newlist{pasoslista}{enumerate}{1}
\setlist[pasoslista]{
  label=\textcolor{milcGradeDark}{\sffamily\bfseries\arabic*.},
  leftmargin=8mm, labelsep=3mm, itemsep=2.8mm,
  topsep=1mm, parsep=0pt, partopsep=0pt, before=\RaggedRight
}

\newcommand{\phasecard}[4]{
\begin{tcolorbox}[enhanced,colback=#3,colframe=#2,arc=3mm,boxrule=.5pt,height=3.05cm,valign=center,halign=center,left=1.5mm,right=1.5mm,top=2mm,bottom=2mm]
{\sffamily\bfseries\fontsize{22}{24}\selectfont\textcolor{#2}{#1}}\par
{\sffamily\bfseries\small #4}
\end{tcolorbox}
}

% ── Piezas del rediseno 2026 (legibilidad 6.º-11.º) ───────────────────
% Banda de estacion: reemplaza el titlebox macizo. Titulo a la izquierda,
% dato practico (tiempo/modalidad) a la derecha, en una sola linea.
% Nunca huerfana: pide 9 lineas libres (o salta de pagina rellenando con
% \vfill) y prohibe el salto justo despues de la banda. Nueve y no seis:
% la caja partible que sigue necesita ~80pt (titulo adosado, relleno y dos
% lineas) para arrancar en la misma pagina; con menos, tcolorbox la manda
% entera a la siguiente y la banda se queda sola. El penalty 10000 va
% en `after`, ANTES del espacio posterior: un `after skip` (aunque sea 0pt)
% es un \vskip tras la caja, y ese glue es un punto de corte legal que un
% \nopagebreak escrito despues de \end{tcolorbox} ya no cierra. Cada banda
% deja marcador PDF y actualiza la cabecera derecha.
\newcounter{milcestacion}
\newcommand{\milcestacion}[2]{%
\Needspace*{9\baselineskip}%
\stepcounter{milcestacion}%
\pdfbookmark[1]{#2}{milcestacion\themilcestacion}%
\markright{#2}%
\begin{tcolorbox}[enhanced,colback=#1,colframe=#1,arc=2mm,boxrule=0pt,
  left=5mm,right=5mm,top=2.6mm,bottom=2.6mm,before skip=11pt,
  after={\par\nopagebreak\vskip7pt}]
{\sffamily\bfseries\fontsize{15}{18}\selectfont\milctintasobre{#1}#2}
\end{tcolorbox}}

% Tarjeta de paso numerada: sustituye las tablas «Pilar / Aplicacion», que
% eran formato de consulta docente, no de estudiante. El numero va en un
% circulo sobre el borde para que la secuencia se lea de un vistazo.
\newcommand{\milcpaso}[3]{%
\Needspace{4\baselineskip}%
\begin{tcolorbox}[enhanced,colback=white,colframe=milcLinea,arc=1.5mm,boxrule=.9pt,
  left=14mm,right=4mm,top=3mm,bottom=3mm,before skip=4pt,after skip=4pt,
  fontupper=\RaggedRight,
  overlay={\node[anchor=north west,circle,fill=#2,text=white,inner sep=0pt,
    minimum size=7.4mm,font=\sffamily\bfseries\small]
    at ([xshift=3.6mm,yshift=-3.2mm]frame.north west) {#1};}]
#3
\end{tcolorbox}}

% Casilla marcable de tamano util con lapiz (la anterior era \fbox{\phantom{x}},
% de alto variable segun el texto que la seguia).
\newcommand{\milcmarca}{\framebox[3.8mm]{\rule{0pt}{3.4mm}}}

% Fila marcable de lista suelta (checks de la estacion 1).
\newcommand{\milccheckrow}[1]{%
\par\vspace{1.8mm}\noindent
\makebox[6.4mm][l]{\raisebox{-.55mm}{\textcolor{milcNegro}{\milcmarca}}}%
\parbox[t]{\dimexpr\linewidth-6.4mm\relax}{\RaggedRight #1}\par}

% Celda marcable dentro de la rubrica (tabla de autoevaluacion). Misma
% sangria francesa, pero pensada para vivir dentro de una columna X.
\newcommand{\milcopcion}[1]{%
\makebox[5.2mm][l]{\raisebox{-.55mm}{\textcolor{milcNegro}{\milcmarca}}}%
\parbox[t]{\dimexpr\linewidth-5.2mm\relax}{\RaggedRight #1}}

% ── Recursos visuales de la guía ──────────────────────────────────────
% Declarados en el bloque `recursos:` del YAML y emitidos por el builder.
% \guiaFigura   → imagen rasterizada o PDF (foto, carta celeste, montaje).
% \guiaDiagrama → fuente TikZ/circuitikz (.tex) incrustada como vector:
%                 es la vía para circuitos y esquemáticos editables.
% [#1] = texto alternativo (alt) · #2 = ruta relativa al .tex · #3 = pie
% El alt llega al PDF etiquetado (/Alt) y lo lee el lector de pantalla.
\newcommand{\guiaFigura}[3][]{%
\begin{center}
\includegraphics[width=0.88\linewidth,height=0.38\textheight,keepaspectratio,alt={#1}]{#2}\\[1.5mm]
{\footnotesize\itshape\textcolor{milcNegro!75}{#3}}
\end{center}
}

\newcommand{\guiaDiagrama}[2]{%
\begin{center}
\input{#1}\\[1.5mm]
{\footnotesize\itshape\textcolor{milcNegro!75}{#2}}
\end{center}
}

\begin{document}
\thispagestyle{empty}

% ── Portada ───────────────────────────────────────────────────────────
% Antes: un rectangulo de color a pagina completa. Se veia bien en pantalla,
% pero en la fotocopiadora del colegio se comia el toner y salia gris sucio.
% Ahora la banda ocupa el tercio superior y el resto va en blanco.
\begin{tikzpicture}[remember picture,overlay]
  \path[fill=milcGradePrimary]
    (current page.north west) rectangle ([yshift=-10.6cm]current page.north east);

  \node[anchor=north west,text=milcGradeText] at ([xshift=2.0cm,yshift=-1.55cm]current page.north west)
    {\sffamily\bfseries\fontsize{12}{14}\selectfont GRADO};
  \node[anchor=north west,text=milcGradeText,opacity=.85] at ([xshift=2.0cm,yshift=-2.35cm]current page.north west)
    {\sffamily\bfseries\fontsize{8.5}{10}\selectfont SERIE GUÍAS · TECNOLOGÍA E INFORMÁTICA};
  \node[anchor=north east,text=milcGradeText,opacity=.85,align=right] at ([xshift=-2.0cm,yshift=-1.55cm]current page.north east)
    {\sffamily\bfseries\fontsize{9}{12}\selectfont I.E. Sor María Juliana\\Cartago, Valle del Cauca};

  \node[anchor=west,text=milcGradeText] (gradonum) at ([xshift=1.85cm,yshift=-7.1cm]current page.north west)
    {\sffamily\bfseries\fontsize{112}{112}\selectfont 8};
  % El circulito de grado (°) solo aplica a las guias de grado («11°»); en
  % Bebras y Semillero el numero grande es un momento/modulo, no un grado.
  % Se posiciona relativo al nodo `gradonum`, no en coordenadas absolutas.
  \draw[milcGradeText,line width=1.6pt]
    ([xshift=2.0mm,yshift=-4.5mm]gradonum.north east) circle (.15cm);

  \node[anchor=west,text=milcGradeText,text width=11.4cm,align=left] at ([xshift=1.5cm,yshift=0.5cm]gradonum.east)
    {
     {\sffamily\bfseries\fontsize{9}{11}\selectfont\MakeUppercase{Período 2 · Lógica, algoritmos y computación física}}\\[3mm]
     {\sffamily\bfseries\fontsize{21}{25}\selectfont\setlength{\baselineskip}{27pt}Umbrales\\[4pt]que se calibran}\\[3.5mm]
     {\sffamily\bfseries\fontsize{15}{18}\selectfont Guía 16-8-TIC}
    };
\end{tikzpicture}

\vspace*{9.4cm}
\vfill

{\sffamily\bfseries\fontsize{8}{10}\selectfont\textcolor{milcGradeDark}{LO QUE TE LLEVAS HOY}}

\begin{tcolorbox}[enhanced,colback=white,colframe=milcNegro,arc=2mm,boxrule=1.6pt,
  left=5mm,right=5mm,top=4mm,bottom=4mm,before skip=3pt,after skip=12pt,
  fontupper=\RaggedRight\fontsize{11}{15.5}\selectfont]
Un programa en MakeCode que lea el sensor de luz o el de temperatura, guarde la lectura en una variable con nombre y dispare tres alertas distintas según tres umbrales. Y una bitácora de calibración con al menos cinco mediciones reales, los tres umbrales decididos y una línea que justifique cada uno.
\end{tcolorbox}

\begin{tcolorbox}[enhanced,colback=milcGris,colframe=milcGris,arc=2mm,boxrule=0pt,
  left=5mm,right=5mm,top=3.5mm,bottom=3.5mm,after skip=12pt]
{\sffamily\bfseries\fontsize{8}{10}\selectfont\textcolor{milcNegro!55}{ESTA GUÍA ES DE}}\\[3mm]
\begin{tabularx}{\linewidth}{X p{3.2cm} p{3.2cm}}
\linead{\linewidth} & \linead{3.0cm} & \linead{3.0cm}\\[-1mm]
{\fontsize{8.5}{10}\selectfont\textcolor{milcNegro!55}{Nombre completo}} &
{\fontsize{8.5}{10}\selectfont\textcolor{milcNegro!55}{Grupo}} &
{\fontsize{8.5}{10}\selectfont\textcolor{milcNegro!55}{Fecha}}\\
\end{tabularx}
\end{tcolorbox}

\vfill

\noindent\textcolor{milcLinea}{\rule{\linewidth}{1pt}}\\[2mm]
{\sffamily\bfseries\fontsize{9.5}{12}\selectfont Autor: Dr. Álvaro Cárdenas Orozco}\\[1mm]
{\sffamily\fontsize{8.5}{11}\selectfont\textcolor{milcNegro!55}{Metodología MILC · apertura ancestral · pensamiento computacional · triángulo Dussel–estoicismo–Floridi}}

\newpage

\section*{Apertura: Variables y umbrales --- calibrar antes de programar}

\begin{softbox}{milcGradeDark}{milcGradeSoft}{Saber ancestral}
En Toribío, en el norte del Cauca, los mayores nasa leen el tiempo en señales que pocos saben interpretar. Si las hormigas arrieras salen al aire libre a recoger comida, va a llover. Si el sapo croa, va a llover. Y la luna avisa con su color: \textbf{nus ate}, la luna de lluvia, se ve pálida y blanca; \textbf{sek ate}, la luna de sol, se ve amarillenta (Ramos García, Tenorio y Muñoz Yule, 2011). Cada señal es un sensor con un \textbf{umbral}: no basta con ver la luna, hay que saber a partir de qué tono anuncia agua. Y ese umbral se vuelve a ajustar. Don Luis Ardo Ascué lo cuenta así: antes, apenas el sapo croaba, caía un torrente; ahora solo cae una llovizna. La señal sigue ahí, pero ya no anuncia lo mismo. La \textbf{cara de exclusión}: los autores del capítulo advierten que no buscan probar si las señales aciertan, sino entender cómo ve el mundo un pueblo. Y Toribío lee esas señales en medio de un conflicto armado que las marca. Hoy vas a hacer lo mismo: medir varias veces antes de decidir el umbral, y volver a mirarlo cuando la realidad cambie.\par\smallskip{\footnotesize\textcolor{milcNegro!70}{\textbf{Fuente.} Ramos García, C., Tenorio, A. D. y Muñoz Yule, F. (2011). Ciclos naturales, ciclos culturales: percepción y conocimientos tradicionales de los nasas frente al cambio climático en Toribío, Cauca, Colombia. En A. Ulloa (Ed.), \emph{Perspectivas culturales del clima} (pp. 247--274). Universidad Nacional de Colombia.}}
\end{softbox}

\begin{softbox}{milcTurquesa}{milcGris}{Saber contemporáneo}
El sensor de luz del micro:bit entrega un número entre 0, oscuridad total, y 255, luz fuerte. El de temperatura entrega grados Celsius. La pregunta ya no es «¿hay luz?», sino «¿cuánta luz hay y desde qué valor disparo la alerta?». Ese valor es el \textbf{umbral}: la frontera entre una zona y otra. Los umbrales no son universales. Un aula con cortinas abiertas puede estar «oscura» en 100; la misma aula con cortinas cerradas, en 40. Por eso los umbrales se \textbf{calibran}: se toman varias lecturas en condiciones reales antes de decidir. Para comparar lecturas necesitas guardarlas. Eso hacen las \textbf{variables}: «establecer nivel\_luz a nivel de luz» guarda el número, y después «si nivel\_luz $<$ umbral\_oscuro» decide. Una variable con nombre claro, nivel\_luz y no x, se entiende un mes después. Sin variables no puedes comparar; sin calibración, comparas contra un número inventado.
\end{softbox}

\begin{softbox}{milcMostaza}{milcGris}{Pregunta-puente}
\textit{Don Luis notó que el sapo ya no anuncia el torrente de antes y ajustó lo que espera de esa señal. Si tu alarma de «falta luz» sonara a mediodía con el sol entrando, ¿qué habría fallado: el sensor o el umbral que tú pusiste?}
\end{softbox}

\begin{softbox}{milcVerde}{milcGris}{Saber y saber hacer}
Al terminar podrás: (1) \textbf{identificar} un umbral razonable a partir de una lista de lecturas; (2) \textbf{aplicar} una calibración de cinco mediciones reales con bitácora; (3) \textbf{evaluar} tres umbrales y programarlos con variables y tres alertas; (4) \textbf{explicar} por qué el mismo umbral no sirve en otro salón.
\end{softbox}



\small
\begin{tabularx}{\textwidth}{>{\bfseries}p{3.0cm}Y}
\rowcolor{milcGradePrimary!12}Autor & Dr. Álvaro Cárdenas Orozco\\
DBA & Pensamiento computacional aplicado a sistemas de control con sensores y actuadores (MEN, T\&I)\\
Referentes & Pensamiento computacional · MakeCode / micro:bit · Logos como razón universal\\
Producto final & Un programa en MakeCode que lea el sensor de luz o el de temperatura, guarde la lectura en una variable con nombre y dispare tres alertas distintas según tres umbrales. Y una bitácora de calibración con al menos cinco mediciones reales, los tres umbrales decididos y una línea que justifique cada uno.\\
\end{tabularx}
\normalsize

\puente{En la sesión 5 el micro:bit se expresó con luces y sonido. Hoy aprende a medir cuánto, no solo si, y a decidir con umbrales que tú calibraste. En la sesión 7 vas a buscar los errores cuando el programa no haga lo que esperas.}

\milcestacion{milcGradeDark}{Ruta MILC: así vamos a aprender}
\begin{tabularx}{\textwidth}{>{\centering\arraybackslash}X>{\centering\arraybackslash}X>{\centering\arraybackslash}X>{\centering\arraybackslash}X}
\phasecard{1}{milcVerde}{milcGris}{Escucha\\[-1mm]\small Observo, escucho y pregunto.} &
\phasecard{2}{milcTurquesa}{milcGris}{Entender\\[-1mm]\small Sistematizo y ordeno.} &
\phasecard{3}{milcMagenta}{milcGris}{Hacer\\[-1mm]\small Construyo y aplico.} &
\phasecard{4}{milcVino}{milcGris}{Cierre\\[-1mm]\small Evalúo y reflexiono.}
\end{tabularx}


\puente{Antes de tocar el micro:bit, vas a decidir un umbral con lápiz sobre cinco lecturas de un aula. Es la lectura de la luna, pero con números.}

\milcestacion{milcVerde}{Estación 1 de 3 · Escucha}

\begin{softbox}{milcVerde}{milcGris}{Escena para escuchar}
\textbf{Actividad 1 · IDENTIFICA --- El umbral sobre cinco lecturas} (15 min · individual). (1) Copia en el cuaderno estas lecturas del sensor de luz de un aula: 7 a. m. = 120; 9:30 a. m. = 180; 12 m. = 210; 2 p. m. = 195; 5 p. m. = 65. (2) Quieres una alarma de «falta luz» que suene solo cuando el aula esté oscura de verdad. Escribe qué pasaría en cada uno de los cinco momentos si el umbral fuera 100. (3) Escribe qué pasaría si el umbral fuera 150. (4) Decide tu umbral y escribe en dos líneas por qué. (5) Anota qué lectura te haría dudar de tu umbral.
\end{softbox}

{\sffamily\bfseries\fontsize{8}{10}\selectfont\textcolor{milcNegro!55}{MARCA CUANDO LO HAYAS HECHO}}
\milccheckrow{Tengo las cinco lecturas copiadas con su hora.}
\milccheckrow{Escribí qué pasa en los cinco momentos con umbral 100 y con umbral 150.}
\milccheckrow{Decidí un umbral y lo justifiqué en dos líneas.}

\begin{infoband}{milcVerde}


\textbf{Lo que va al cuaderno (Actividad 1 · IDENTIFICA).} \textbf{Título:} «Actividad 1 --- El umbral sobre cinco lecturas». \textbf{Formato:} las cinco lecturas con su hora, dos filas de «qué pasa» (umbral 100 y umbral 150) y tu umbral con su justificación. \textbf{Extensión:} media página. \textbf{Sabes que terminaste cuando} puedes decir en voz alta en cuáles de los cinco momentos sonaría tu alarma y por qué en esos.
\end{infoband}

\puente{Ya decidiste un umbral en papel y viste qué pasaba si lo movías. Ahora aprendes las variables y, con tu pareja, haces cinco mediciones reales en el salón.}

\milcestacion{milcTurquesa}{Estación 2 de 3 · Entender}

\begin{softbox}{milcTurquesa}{milcGris}{Lo que vas a entender}
\textbf{Actividad 2 · APLICA --- Cinco mediciones reales} (30 min · parejas). (1) Con tu pareja, en MakeCode, arrastren «para siempre» y adentro «mostrar número (nivel de luz)»; carguen al micro:bit o usen el simulador. (2) Dibujen en el cuaderno la bitácora de cuatro columnas: momento, lectura, condición real, observación. (3) Tomen cinco lecturas en condiciones distintas: micro:bit sobre la mesa, junto a la ventana, bajo la lámpara, tapado con la mano, dentro del bolso. (4) Anoten cada lectura en su fila con lo que observaron. (5) Ordenen las cinco lecturas de menor a mayor y marquen cuáles son «normales».
\end{softbox}

\milcpaso{1}{milcTurquesa}{\textbf{La variable guarda la lectura.} En la categoría Variables, «establecer nivel\_luz a (nivel de luz)» guarda el número que dio el sensor en ese instante. «Cambiar nivel\_luz por 1» le suma uno. El nombre lo eliges tú: nivel\_luz, temperatura\_aula, umbral\_oscuro. Con x o y, en una semana no sabes qué guardaste.}
\milcpaso{2}{milcTurquesa}{\textbf{La bitácora de calibración.} Cuatro columnas: momento (hora o condición), lectura (el número exacto), condición real (¿estaba oscuro?, ¿había sol?) y observación (la cortina, la lámpara, la mano). Cinco mediciones en condiciones distintas son el mínimo para decidir un umbral.}
\milcpaso{3}{milcTurquesa}{\textbf{Cómo se decide un umbral.} Ordena las lecturas de menor a mayor. Mira dónde se juntan la mayoría: ese es el rango normal. El umbral inferior va por debajo de la menor lectura normal; el superior, por encima de la mayor. Con 65, 120, 180, 195 y 210, el rango normal es 120 a 210. Oscuro podría ser menos de 100 y brillante más de 220.}
\milcpaso{4}{milcTurquesa}{\textbf{Cómo se hace, paso a paso.} (1) Mostrar la lectura en pantalla. (2) Dibujar la bitácora. (3) Medir cinco veces en condiciones distintas. (4) Anotar cada fila. (5) Ordenar y marcar el rango normal.}

\begin{drawbox}{milcTurquesa}{La calibración, en una mirada}
\textbf{Lectura:} el número que da el sensor. \\[1mm] \textbf{Variable:} donde se guarda la lectura, con nombre claro. \\[1mm] \textbf{Bitácora:} momento, lectura, condición, observación. \\[1mm] \textbf{Rango normal:} donde se juntan la mayoría de lecturas. \\[1mm] \textbf{Umbral:} la frontera, por fuera del rango normal.
\end{drawbox}

\begin{softbox}{milcMostaza}{milcGris}{Errores comunes y su corrección}
(1) \textbf{Umbral inventado}: la alarma suena a mediodía o no suena de noche. (2) \textbf{Variables x, y, z}: nadie sabe qué guardan. (3) \textbf{Una o dos mediciones}: no alcanzan para ver el rango normal. (4) \textbf{Umbral pegado al rango normal}: alarma falsa cada rato. (5) \textbf{No volver a mirar}: el umbral falló en la prueba y quedó igual.
\end{softbox}

\begin{infoband}{milcTurquesa}


\textbf{Lo que va al cuaderno (Actividad 2 · APLICA).} \textbf{Título:} «Actividad 2 --- Cinco mediciones reales». \textbf{Formato:} la bitácora de cuatro columnas con cinco filas llenas, y debajo las lecturas ordenadas con el rango normal marcado. \textbf{Extensión:} media página. \textbf{Sabes que terminaste cuando} las cinco filas tienen lectura y observación, y puedes señalar cuál es la lectura normal más baja y la más alta.
\end{infoband}

\puente{Ya tienes la bitácora con cinco lecturas. Toca decidir tres umbrales, programarlos con variables y comprobar que las tres alertas suenan cuando deben.}

\milcestacion{milcMagenta}{Estación 3 de 3 · Hacer}

\begin{softbox}{milcMagenta}{milcGris}{Cómo vas a construir tu producto}
\textbf{Actividad 3 · EVALÚA --- Tres umbrales, tres alertas} (25 min · individual). (1) Con tu bitácora, decide tres zonas: oscuro, normal y brillante, o frío, ambiente y caliente. Escribe los dos umbrales que las separan y una línea de justificación por cada uno. (2) En MakeCode, dentro de «para siempre», pon «establecer nivel\_luz a (nivel de luz)». (3) Agrega un «si… si no, si… si no» con tres ramas: nivel\_luz $<$ umbral\_oscuro, nivel\_luz $>$ umbral\_brillante y el resto. (4) Pon en cada rama un ícono distinto: luna, sol y una equis pequeña. (5) Prueba tapando el micro:bit y poniéndolo bajo la luz. Si una alerta sale en el momento equivocado, mueve el umbral y anota el cambio.
\end{softbox}

\milcpaso{1}{milcMagenta}{\textbf{Tu entrega tiene tres piezas.} (1) El programa con la variable y las tres ramas, compartido con enlace o archivo .hex. (2) La bitácora con las cinco mediciones y los tres umbrales justificados. (3) La nota del ajuste, si tuviste que mover un umbral después de probar.}
\milcpaso{2}{milcMagenta}{\textbf{Calibrar primero, programar después.} El programa se escribe en cinco minutos cuando los umbrales ya están decididos. Lo que toma tiempo es medir. Quien programa primero termina cambiando números al azar hasta que «funciona» una vez.}
\milcpaso{3}{milcMagenta}{\textbf{Tres alertas que se distingan.} Luna para oscuro, sol para brillante, equis para normal. O tres notas: grave, media, aguda. Si dos alertas se parecen, quien mira el micro:bit no sabe qué pasa.}
\milcpaso{4}{milcMagenta}{\textbf{La prueba y el ajuste.} Tapas el micro:bit: debe salir la luna. Lo pones bajo la lámpara: debe salir el sol. Si sale la equis en la oscuridad, tu umbral de oscuro quedó muy bajo. Súbelo y anota por qué. Ese ajuste es lo que hizo don Luis con el sapo.}

\begin{drawbox}{milcMagenta}{Antes de entregar}
\checkbox Cinco mediciones reales en condiciones distintas. \\[1mm] \checkbox Bitácora con las cuatro columnas llenas. \\[1mm] \checkbox Dos umbrales decididos, con una línea de justificación cada uno. \\[1mm] \checkbox Variable con nombre claro, no x ni y. \\[1mm] \checkbox Tres ramas que muestran tres alertas distintas. \\[1mm] \checkbox Probado tapando y bajo la luz; ajuste anotado si hizo falta. \\[1mm] \checkbox Programa compartido con enlace o archivo .hex.
\end{drawbox}

\begin{softbox}{milcMagenta}{milcGris}{Plantilla de trabajo}
\textbf{Sensor.} \textit{Luz o temperatura.} \\[1mm] \textbf{Lecturas.} \textit{Las cinco, de menor a mayor.} \\[1mm] \textbf{Rango normal.} \textit{De … a …} \\[1mm] \textbf{Umbrales.} \textit{Oscuro si menos de …, porque …; brillante si más de …, porque …} \\[1mm] \textbf{Ajuste.} \textit{Moví … porque en la prueba …}
\end{softbox}

\begin{infoband}{milcMagenta}


\textbf{Lo que va al cuaderno (Actividad 3 · EVALÚA).} \textbf{Título:} «Actividad 3 --- Tres umbrales, tres alertas». \textbf{Formato:} los dos umbrales con su justificación de una línea, el dibujo de las tres ramas con su ícono, y la nota del ajuste si lo hubo. \textbf{Extensión:} media página. \textbf{Sabes que terminaste cuando} al tapar el micro:bit sale la luna, bajo la lámpara sale el sol, y cada umbral tiene su porqué escrito.
\end{infoband}

\puente{Lo que entregas hoy: el programa con variable y tres alertas, y la bitácora de calibración con los umbrales justificados. La prueba de calidad: cubres el micro:bit con la mano y sale la alerta de oscuro; lo pones bajo la lámpara y sale la de brillante.}

\milcestacion{milcMostaza}{Producto de la sesión}
\textbf{Programa con variable y tres alertas por umbral + bitácora de calibración con cinco mediciones y umbrales justificados}.
\begin{softbox}{milcMostaza}{milcGris}{Criterios de calidad}
(1) La bitácora tiene cinco mediciones reales en condiciones distintas, con las cuatro columnas. (2) Los dos umbrales están por fuera del rango normal y cada uno tiene una justificación. (3) El programa guarda la lectura en una variable con nombre claro. (4) Las tres ramas muestran tres alertas distintas. (5) Al tapar el micro:bit sale la alerta de oscuro y bajo la luz la de brillante. (6) El programa está compartido con enlace o archivo .hex.
\end{softbox}

\puente{Antes de cerrar, mira lo que hiciste desde cinco lados. Un umbral que salió de cinco lecturas vale más que uno que se te ocurrió.}

\milcestacion{milcVino}{Evaluación liberadora · cinco dimensiones}

\begin{softbox}{milcVino}{milcGris}{Sentido de la evaluación}
Evaluar no es solo revisar una nota. Es reconocer qué aprendí, qué cambió en mí, qué emociones manejé, cómo me relaciono con la ciudad y la comunidad, y qué le dejaré a quien viene detrás.
\end{softbox}

\textbf{Mi reflexión liberadora en cinco dimensiones.} Completa cada frase iniciada:

\small
\begin{tabularx}{\textwidth}{>{\bfseries\RaggedRight\arraybackslash}p{3.4cm}Y>{\RaggedRight\arraybackslash}p{4.6cm}}
\rowcolor{milcVino}\color{white}Dimensión & \color{white}\bfseries Mi respuesta (frame starter) & \color{white}\bfseries Algo que descubrí\\
1. Personal & Lo que más cambió en mí fue \linead{6.5cm} & \linead{4.2cm}\\
2. Emocional & La emoción más fuerte fue \linead{2.5cm} y la regulé \linead{3cm} & \linead{4.2cm}\\
3. Ciudadana & En democracia esto importa porque \linead{6cm} & \linead{4.2cm}\\
4. Local (Cartago) & En mi barrio sería útil para \linead{6.5cm} & \linead{4.2cm}\\
5. Intergeneracional & A mi abuela/abuelo le diría \linead{6.5cm} & \linead{4.2cm}\\
\end{tabularx}
\normalsize

\begin{infoband}{milcVino}
\textbf{Para la próxima sustentación.} Vuelve a esta tabla en seis meses. Verás que las respuestas cambian — no porque lo aprendido cambie, sino porque tú cambias. La evaluación liberadora es una conversación contigo a través del tiempo.
\end{infoband}


\puente{Para terminar, tres ideas sobre medir antes de decidir, contadas con palabras de esta guía: Dussel sobre la información que nunca se lee sola, Epicteto sobre las cosas y las opiniones, y la Onlife Initiative sobre las alarmas que gastan la atención.}

\milcestacion{milcGradeDark}{Tres ideas para cerrar · Dussel, un estoico y Floridi}

\begin{softbox}{milcVino}{milcGris}{Enrique Dussel · lente del nosotros}
\textit{Es ingenuo creer que la información se lee sola, sin conflictos ni contexto.}

{\footnotesize\itshape\color{milcNegro!70}--- idea tomada de Enrique Dussel, contada con palabras de esta guía. No es una cita textual.}

El número 120 no dice «oscuro» por sí mismo. Lo dice el salón, la hora y la cortina. Un umbral es una decisión sobre el contexto, y por eso el de tu salón no sirve en otro.

\textbf{¿Qué número leí esta semana como si hablara solo, sin preguntar de dónde salió?}
\end{softbox}
\linead{\linewidth}\\[1mm]
\linead{\linewidth}

\begin{softbox}{milcMostaza}{milcGris}{Estoicismo · Epicteto · cuidado interior}
\textit{No nos inquietan las cosas, sino las opiniones que tenemos sobre ellas.}

{\footnotesize\itshape\color{milcNegro!70}--- idea tomada de Epicteto, contada con palabras de esta guía. No es una cita textual.}

La lectura del sensor es la cosa. El umbral que inventaste sin medir es una opinión. Cuando la alarma suena a mediodía, no falló el sensor: falló la opinión.

\textbf{¿Cuál de mis umbrales de hoy era una opinión antes de medir?}
\end{softbox}
\linead{\linewidth}\\[1mm]
\linead{\linewidth}

\begin{softbox}{milcTurquesa}{milcGris}{Luciano Floridi y la Onlife Initiative · lente de la infoesfera}
\textit{El diseño de nuestras tecnologías, empezando por sus valores por defecto, debería proteger la capacidad de atención de las personas.}

{\footnotesize\itshape\color{milcNegro!70}--- idea tomada de Luciano Floridi y la Onlife Initiative, contada con palabras de esta guía. No es una cita textual.}

Una alarma que suena cuando no toca enseña a ignorarla. Calibrar el umbral es respetar la atención de quien va a oírla: suena poco y suena cuando importa.

\textbf{¿Qué alerta de mi celular suena tanto que ya no la miro?}
\end{softbox}
\linead{\linewidth}\\[1mm]
\linead{\linewidth}

\begin{infoband}{milcNegro}
\textbf{Nota para el docente.} Las tres ideas están contadas con palabras de esta guía a partir del pensamiento de cada autor; no son citas textuales. De 9.º en adelante se trabajan con la obra y su referencia.
\end{infoband}

\milcestacion{milcVino}{Mi autoevaluación · marca una casilla por fila}

{\small
\setlength{\extrarowheight}{2.2pt}
\begin{tabularx}{\textwidth}{>{\RaggedRight\arraybackslash\bfseries}p{3.0cm}YYY}
\rowcolor{milcVino}\color{white}\bfseries Criterio & \color{white}\bfseries Lo logré & \color{white}\bfseries En proceso & \color{white}\bfseries Necesito apoyo\\[.6mm]
Comprensión del tema & \milcopcion{Explico las ideas con mis palabras.} & \milcopcion{Reconozco las ideas, falta precisión.} & \milcopcion{Necesito repasar lo básico.}\\[.8mm]
\rowcolor{milcGris}Pensamiento computacional & \milcopcion{Usé los pilares al armar mi producto.} & \milcopcion{Usé algunos pilares.} & \milcopcion{Necesito releer la estación 2.}\\[.8mm]
Aplicación y producto & \milcopcion{Mi producto cumple los criterios.} & \milcopcion{Está armado, con ajustes pendientes.} & \milcopcion{Aún está en construcción.}\\[.8mm]
\rowcolor{milcGris}Reflexión 5 dimensiones & \milcopcion{Respondí cada una con honestidad.} & \milcopcion{Respondí algunas con detalle.} & \milcopcion{Mis respuestas son muy generales.}\\[.8mm]
Triángulo de pensamiento & \milcopcion{Conecté las 3 voces con mi trabajo.} & \milcopcion{Conecté al menos una.} & \milcopcion{No las relaciono con mi proceso.}\\[.8mm]
\end{tabularx}}

\vspace{3mm}
\textbf{Hoy entendí que:}\\[1mm]
\linead{\linewidth}\\[1mm]
\linead{\linewidth}

\textbf{Mi compromiso para la próxima sesión es:}\\[1mm]
\linead{\linewidth}\\[1mm]
\linead{\linewidth}

\begin{softbox}{milcMostaza}{milcGris}{Cita de despedida · Marco Aurelio}
\textit{``El obstáculo en el camino se vuelve el camino. Lo que estorba al camino se vuelve camino.''}\\[1mm]
Cada error de hoy, bien observado, se convierte en pieza del camino que pisarás mañana.
\end{softbox}

\small
\begin{tabularx}{\textwidth}{>{\bfseries}p{3.6cm}Y}
\rowcolor{milcGradeDark!12}Nombre completo: & \linead{10cm}\\
Fecha: & \linead{4cm} \quad Curso/grupo: \linead{4cm}\\
Mi firma: & \linead{9cm}\\
\end{tabularx}
\normalsize

\begin{infoband}{milcGradeDark}
\textbf{Para la próxima sesión.} Llega con tu programa calibrado y la bitácora. En la sesión 7 vas a buscar el error cuando el programa no hace lo que esperas: una hipótesis, una prueba. Los umbrales que calibraste hoy son el primer lugar donde vas a mirar.
\end{infoband}

\end{document}
